%/*% pdflatex -file-line-error -interaction nonstopmode %
% Comment
\documentclass{article}
\usepackage{amsmath}
\usepackage{amsfonts}
\usepackage{hyperref}
\usepackage{exercises}
\usepackage{amsthm}
\usepackage{yhmath}
\newcommand*{\oldneg}{\mathord{\sim}}

\theoremstyle{definition}
\newtheorem{definition}{Definition}

\setlength{\parindent}{0pt}
\setlength{\parskip}{10pt}

\title{Chapter 13 Answers}
\begin{document}
\maketitle

\section*{Problem 1}
(a): Let us note that, because the Group $I$ axioms are just the system $S_1$ of Chapter $7$,
it follows that all tautologies are provable in Peano Arithmetic (as 
$S_1$ was proved to be complete). Then, pick any formula $C$ such that $x$ has no free occurrences (such as a propositional formula). Then, it follows 
that $C\supset F(x)$ is provable (as this is a tautology if $F(x)$ is provable), and therefore so is $\forall xF(x)$ by $Group\ Inference\ Rules\ Rule\ 2$.
\newline\newline 
(b): If $F(x)$ is provable, then $\forall xF(x)$ is provable (by (a)). Hence, by $A10$, because $\forall xF(x)\supset F(t)$ is provable, so is 
$F(\overline n)$ for every number $n$ by M.P. (because every number $n$ is a term).\newline\newline
(c): If $F(x,y)$ is provable, then $F(\overline n,y)$ is provable for every $n$ (by (b)), and then $F(\overline n, \overline m)$ is provable for every 
$n$ and $m$ (by (b)).
\section*{Problem 2}
(1):\newline For this, we first introduce into our alphabet $K$ the symbol $v$. Then, we have the following axiom schemes with the predicate $Stx$ 
(for string)
 representing all the strings of $v$'s:\newline $Stv$\newline $Stx\rightarrow Stxv$\newline
(2):\newline Now we introduce the symbols $($ and $)$ into $K$, and then $Vx$ represents all the variables in Peano Arithmetic:\newline
$Stx\rightarrow V(x)$\newline
(3):\newline We need the following axioms for the predicate Dist (for distinct):\newline $stx\rightarrow sty\rightarrow Dist(x)(xy)$\newline
$Distx,y\rightarrow disty,x$\newline
(4):\newline Now we introduce the symbols $0$ and $'$. Then, $Num$ (for numeral) represents the set of Peano Numerals:\newline
$Num0$\newline $Numx\rightarrow Numx'$\newline
(5):\newline We now add the symbols $+$ and $\times$ to $K$. Then $Tx$ (for term) represents the set of all terms:\newline
$Numx\rightarrow Tx$\newline $Vx\rightarrow Tx$\newline $Tx\rightarrow Ty\rightarrow T(x+y)$\newline $Tx\rightarrow Ty\rightarrow T(x\times y)$\newline
$Tx\rightarrow Tx'$\newline
(6):\newline We now add the symbols $=$ and $\le$ to $K$. Then $Afx$ (for atomic formula) represents the set of all atomic formulas:\newline
$Tx\rightarrow Ty\rightarrow Afx=y$\newline $Tx\rightarrow Ty\rightarrow Afx\le y$\newline
(7):\newline We now add the symbols $\oldneg$, $\wedge$, $\vee$, $\supset$, $\forall$ and $\exists$ into $K$. Then $Fx$ represents the set of all formulas:
\newline $Afx\rightarrow Fx$\newline $Fx\rightarrow F\oldneg x$\newline $Fx\rightarrow Fy\rightarrow F(x\wedge y)$\newline $Fx\rightarrow Fy\rightarrow 
F(x\vee y)$\newline $Fx\rightarrow Fy\rightarrow F(x\supset y)$\newline $Fw\rightarrow varx\rightarrow F\forall xw$\newline
$Fw\rightarrow varx\rightarrow f\exists xw$:\newline
(8):\newline We let $Tsub$ be the predicate representing the realtion $t$ is a term, $x$ is a variable, $y$ is a term, and $z$ is the result of 
substituting $t$ for all occurences of $x$ in $y$.\newline $Tt\rightarrow Vx\rightarrow Num y\rightarrow Tsubt,x,y,y$\newline
$Tt\rightarrow Vx\rightarrow Tsubt,x,x,t$\newline $Tt\rightarrow Dist x,y\rightarrow Tsubt,x,y,y$\newline
$Tsubt,x,y,z\rightarrow Tsubt,x,y',z'$\newline
$Tsubt,x,w,z\rightarrow Tsubt,x,v,l\rightarrow Tsubt,x,(w\times v),(z\times l)$\newline
$Tsubt,x,w,z\rightarrow Tsubt,x,v,l\rightarrow Tsubt,x(w+v),(z+l)$\newline
$Note:$ Smullyans answer omits parenthesis (where I placed them), and this seeems to be incorrect, as he defined a term (in those cases) with parenthesis.
\newline
(9):\newline We let $Tfsub$ be the predicate representing the relation $t$ is a term, $x$ is a variable, $y$ is a formula, and $z$ is the result of 
substituting $t$ for all free occurences of $x$ in $y$:\newline
$Tsubt,x,y,z\rightarrow Tsubt,x,w,v\rightarrow Tfsubt,x,y=w,z=v$\newline $Tsubt,x,y,z\rightarrow Tsubt,x,w,v\rightarrow Tfsubt,x,y\le w,z\le v$\newline
$Tfsubt,x,y,z\rightarrow Tfsubt,x,w,v\rightarrow Tfsubt,x,y\wedge w,z\wedge v$\newline $Tfsubt,x,y,z\rightarrow Tfsubt,x,w,v\rightarrow Tfsubt,x,y
\vee w,z\vee v$\newline $Tfsubt,x,y,z\rightarrow Tfsubt,x,w,v\rightarrow Tfsubt,x,y\supset w,z\supset v$\newline $Tfsubt,x,y,z\rightarrow Tfsubt,x,\oldneg y
\oldneg z$\newline $Fy\rightarrow \rightarrow Tt\rightarrow Vx\rightarrow Tfsubt,x,\forall xy,\forall xy$\newline 
$Fy\rightarrow Tt\rightarrow Vx\rightarrow Tfsubt,x,\exists xy,\exists xy$\newline
$Distx,y\rightarrow Tfsubt,x,w,v\rightarrow Tfsubt,x,\forall yw,\forall yv$\newline $Distx,y\rightarrow Tfsubt,x,w,v\rightarrow Tfsubt,x\forall yw,\forall 
yv$\newline
(10):\newline We let $Nocf$ represent the relation $x$ is a variable, $f$ is a formula, and $x$ does not occur free in $f$:\newline
$Distx,y\rightarrow Tsubfy,x,f,f\rightarrow Nocfx,f$\newline
(11):\newline We let $Nb$ represent the relation $x$ and $y$ are variables, $f$ is a formula, and $x$ is not bound by $y$ in $f$:\newline
$Vx\rightarrow Vy\rightarrow Afw\rightarrow Nbx,y,w$\newline $Nbx,y,f_1\rightarrow Nbx,y,f_2\rightarrow Nbx,y,f_1\wedge f_2$\newline
$Nbx,y,f_1\rightarrow Nbx,y,f_2\rightarrow Nbx,y,f_1\vee f_2$\newline $Nbx,y,f_1\rightarrow Nbx,y,f_2\rightarrow Nbx,y,f_1\supset f_2$\newline
$Nbx,y,f\rightarrow Nbx,y,\oldneg f$\newline $Nbx,y,f\rightarrow Nocfx,f\rightarrow Nbx,x,f\rightarrow Nbx,y,\forall yf$\newline 
$Nbx,y,f\rightarrow Nocfx,f\rightarrow Nbx,x,f\rightarrow Nbx,y,\exists yf$\newline $Nbx,y,f\rightarrow Distz,y\rightarrow Nbx,y,\forall zf$\newline
$Nbx,y,f\rightarrow Distz,y\rightarrow Nbx,y,\exists zf$\newline
$Note:$ Smullyan's answer seems to be incorrect, as it allows for $Nbx,y,\forall y\forall xF(x)$ where $Nbx,y,\forall xF(x)$. However, this seems to infact 
be a formula in which $x$ is bounded to $y$, as there is some $G$ (namely $\forall xF(x)$) such that $\forall yG$ is part of $F$ and $x$ occurs in $G$. The 
only time quantifying $y$ will still make $x$ not bound to $y$ is when $x$ does not occur in the formula, and this is why we must have the added condition 
that $x$ is not bound to $x$ (and thus, because it does not occur freely, it cannot occur at all).\newline
(12):\newline We let $Ocf$ represent $t$ is a term, $x$ is a variable, $f$ is a formula, and $t$ is free for $x$ in $f$:\newline
$Nbx,y,f\rightarrow Ocfy,x,f$\newline $Nt\rightarrow Vx\rightarrow Ff\rightarrow Ocft,x,f$\newline $Ocft,x,f\rightarrow Ocft',x,f$\newline 
$Ocft_1,x,f\rightarrow Ocft_2,x,f\rightarrow Ocft_1+t_2,x,f$\newline $Ocft_1,x,f\rightarrow Ocft_2,x,f\rightarrow Ocft_1\times t_2,x,f$\newline
(13):\newline We let the predicate $Ax$ represent mean $x$ is an axiom.\newline $Group\ I\ Axiom\ Schemes:$\newline
$Ff\rightarrow Fg\rightarrow A(f\wedge g)\supset f$\newline $Ff\rightarrow Fg\rightarrow A(f\wedge g)\supset g$\newline $Ff\rightarrow Fg\rightarrow 
Fh\rightarrow A[(f\wedge g)\supset h]\supset [f\supset (g\supset h)]$\newline $Ff\rightarrow Fg\rightarrow Fh\rightarrow A[(f\supset g)\wedge (f\supset
 (g\supset h))]\supset (f\supset h)$\newline $Ff\rightarrow Fg\rightarrow Af\supset (f\vee g)$\newline $Fg\rightarrow Ff\rightarrow Ag\supset (f\vee g)$
\newline $Ff\rightarrow Fg\rightarrow Fh\rightarrow A[(f\supset h)\wedge (g\supset h)]\supset ((f\vee g)\supset h)$\newline
$Ff\rightarrow Fg\rightarrow A((f\supset g)\wedge (f\supset \oldneg g))\supset \oldneg f$\newline $Ff\rightarrow A\oldneg\oldneg f\supset f$\newline
$Group\ II\ Axiom\ Schemes:$\newline $Ocft,x,f\rightarrow Tfsubt,x,f,y\rightarrow A\forall x(f)\supset y$\newline
$Ocft,x,f\rightarrow Tfsubt,x,f,y\rightarrow Ay\supset\exists x(f)$\newline $Group\ III\ Axiom\ Schemes:$\newline
$Tx\rightarrow Ty\rightarrow A(x'=y')\supset (x=y)$\newline $Tx\rightarrow A\oldneg (x'=0)$\newline $Tx\rightarrow Ty\rightarrow Tz\rightarrow
 A((x=y)\wedge (x=z))\supset (y=z)$\newline $Tx\rightarrow Ty\rightarrow A(x=y)\supset (x'=y')$\newline $Tx\rightarrow Ax+0=x$\newline $Tx\rightarrow Ty
\rightarrow Ax+y'=(x+y)'$\newline $Tx\rightarrow Ax\times 0=0$\newline $Tx\rightarrow Ty\rightarrow Ax\times y'=(x\times y)+x$\newline
$Tfsub0,x,f,z\rightarrow Tfsubx',x,f,y\rightarrow A(z\wedge \forall x(f\supset y))\supset \forall x(f)$\newline
$Tx\rightarrow A(x\le 0)\equiv (x=0)$\newline $Tx\rightarrow Ty\rightarrow A(x\le y)\equiv ((x\le y)\vee (x\le y'))$\newline
$Tx\rightarrow Ty\rightarrow A(x\le y)\vee (y\le x)$\newline $Tx\rightarrow Ax=x$\newline $Nn\rightarrow Tfsubn,x,f,y\rightarrow Ay\supset ((x=n)\supset f))$
\newline $Ff\rightarrow A\forall x\oldneg (f)\supset \oldneg\exists x(f)$\newline
$Note:$\newline Smullyans addition of $Vx$ in the mathematical induction scheme is unecessary, as $Tfsubt,x,f,z$ implies that $x$ is a variable.
\newline
(14):\newline We let $Px$ represent the set of provable formulas:\newline $Af\rightarrow Pf$\newline $Px\rightarrow Px\supset y\rightarrow Py$\newline
$Nocfx,c\rightarrow Pc\supset f\rightarrow Pc\supset \forall x(f)$\newline $Nocfx,c\rightarrow Pf\supset c\rightarrow P\exists x(f)\supset c$\newline
$Note:$\newline Smullyans addition of $Fc$ in the last two schemes is unecessary, as this is implied by $Nocfx,c$.\newline
\section*{Problem 3}
(a):\newline We have already shown that $F$ being a formula, $x$ does not occur free in $F$, and $x$ is a variable are all formally representable, and 
therefore the sets of the corresponding Gödel numbers are $\Sigma_1$. Hence, the following formula expresses $S_0$:\newline
$F_0(y)\wedge (\forall x\le y)(V_0(x)\supset Nocf_0(x,y))$\newline
Well, this is equivalent to the following, with $\exists yR_1(x,y)$ expressing $V_0$, $\exists kR_2(w,v,k)$ expressing $Nocf_0(w,v)$ and $\exists lR_3(f,l)$
expressing $F_0(f)$, with all $R_i$ being $\Sigma_0$
:\newline 
$\exists l\exists y\exists v(R_3(f,l)\wedge (\forall z\le f)(R_1(x,y)\supset R_2(z,f,v))))$\newline
which is $\Sigma_1$ by Problem 10(b) of Chapter 11 (because the above formula without the $\exists l$ is $\Sigma_1$ by Problem 10(a) of Chapter 11).\newline 
(b):\newline $x\in P_0$ iff $Pf_0(x)\wedge S_0(x)$\newline and this is $\Sigma_1$ by Problem 13 of Chapter 11.\newline
(c):\newline $x\in R_0$ iff $Pf_0(nx)$ where $n$ is the Gödel number of the negation.\newline 
\section*{Problem 4}
$\Omega_1$:\newline We assume that $m+n=k$. We are to show that $\overline m+\overline n=k$ is provable in Peano Arithmetic. By $A16$ and Proposition 1(b),  
$\overline m+0=\overline m$ is provable. Then, we have $(\overline m+0)'=\overline m'$ being provable by $A15$ and M.P. Then, we also have $\overline m+0'=
(\overline m+0)'$ being provable by A17 and Proposition 1(c), 
and therefore, by A14, we have $\overline m+0'=m'$ being provable. We can repeat this process $n$ times, until we 
eventually end up with $\overline m+\overline n=\overline k$, the desired result.\newline 
$\Omega_2$:\newline We assume that $m\times n=k$. We are to show that $\overline m\times\overline n=\overline k$ is provable in Peano Arithmetic.  
By $A18$ and Proposition 1(b), $\overline m\times 0=0$ is 
provable. Then, by $A19$ and Proposition 1(c), 
$\overline m\times 0'=(\overline m\times 0)+\overline m$ is provable. We repeat the process of applying $A19$ $n$ times, until we 
have proven $\overline m\times \overline n=\overline{m\times n}$\newline
$\Omega_3$:\newline We assume that $m\ne n$, and are to show that $\oldneg (\overline m=\overline n)$ is provable in Peano Arithmetic. Let $m$ be the number 
which is the greater of the two ($m$ and $n$). Then $m-n=k$ for some $k>0$.  
Hence, by $A13$, $\oldneg (\overline k=0)$ is provable. 
By $A12$, $(\overline k'=0')\supset (\overline k=0)$ is provable. Hence, by propositional logic (using the fact 
that $\oldneg X\supset ((Y\supset X)\supset \oldneg Y)$ 
is provable) and M.P. we can prove that $\oldneg (\overline k'=0')$. We can repeat this process $n$ times, 
eventually ending with $\oldneg (\overline m=\overline n)$.\newline
$\Omega_4$:\newline In the base case, $(x\le 0)\equiv (x=0)$ is provable by $A21$. Now assume that for some $n$, $(x\le n)\equiv (x=0\vee ...\vee x=\overline 
n)$ is provable. By $A22$, $(x\le \overline n')\equiv ((x\le n)\vee (x=n'))$. Then, because $(x\le n)\equiv (x=0\vee ...\vee x=\overline n)$ is provable, 
$(x\le \overline n') \equiv (x=0\vee ...\vee x=\overline n\vee x=\overline n')$ (because the formulae $X\equiv Y$ and $Z\equiv (Y\vee W)$ imply the formula 
$Z\equiv (X\vee W)$ by propositional logic), thus completing the induction.\newline
$\Omega_5$:\newline This is immediate from $A23$ and Proposition 1(b).\newline
\section*{Problem 5} $Note:$ From this problem forward, it seems that Smullyan labels the results as general answers for any extension of $(R)$, but 
in fact they rely on arithmetic axioms of PA.\newline 
 We are to show that every true atomic $\Sigma_0$ sentence is 
provable in PA. Let us recall that an atomic $\Sigma_0$ 
formula is either of the form $t_1=t_2, t_1\le t_2$ where $t_1, t_2$ are both terms. In the first case, we can prove this by mathematical induction 
on the number of occurences of plus, times, and the successor operation in the term (it is in fact a double induction). First, we show that for every term 
$t_1=\overline n$ where $t_1=n$. The base case where $t_1$ is a numeral 
is immediate from $A24$. Now, make an induction assumption on the number of plus, 
times, and successor symbols occuring in $t$. So, $t=\overline n$ is provable for some $n$.
By $\Omega_1$, $\overline n+\overline k=\overline m$ for every
 natural number $k$, and $\Omega_2$ proves
 that $\overline n\times \overline k=\overline m$ for every  
natural number $k$, where $m$ is the number that $t+\overline k, t\times \overline k$ designate, respectively,
and the successor operation can be proven just by taking $k$ to 
be $1$ in the addition case. Then, becuase $t=\overline n$ is provable, we can use $A25$ to prove $t+\overline k=\overline m$ or $t\times \overline k=
\overline m$, or $t'=\overline m$.
So, we have shown that, for any term $t$ and any natural number $m$, if $t=m$, then $t=\overline m$ is provable in Peano 
Arithmetic. Next, we must show that $t_1=t_2$ for every terms $t_1, t_2$ such that $t_1=t_2$. Well, by definition, $t_1=t_2$ iff both $t_1,t_2$ designate 
the same natural number $n$, and hence, by the above result, both $t_1=\overline n$ and $t_2=\overline n$ are provable. 
Now, by $A14$, $((x=y)\wedge (x=x))\supset (y=x)$ is provable, and thus, if $x=y$ is provable, so is $y=x$ (using $A24$). Hence, we have 
$((\overline n=t_1)\wedge (\overline n=t_2))$, and thus, by $A14$, $t_1=t_2$ is provable. Now, we need to show that $t_1\le t_2$ is provable for 
every terms $t_1, t_2$, if $t_1\le t_2$ is true. Well, assume that $t_1\le t_2$. Now, because these are terms, by the above, 
$t_1=\overline n, t_2=\overline m$ are provable 
for some numbers $n$ and $m$. There are then two cases: either $n=m$, or one of $m, n$ is greater than the other. In the first case, $t_1\le m$ is provable, 
as $t_1=m$ can be provable after we apply $A14$ and the fact that $m=n$ is provable, and then we use $\Omega_4$. After this, we can use $A25$ to conclude 
that $t_1\le t_2$.
 Now, in the second case, assume (without 
loss of generality) that $m$ is the greater numeral. Well, by $\Omega_4$, $t_1\le \overline m
\equiv (t_1=\overline m\vee...t_1=\overline n\vee...\vee t_1=\overline m$. Because $t_1=\overline n$ is provable, so is $t_1\le \overline m$ (by 
Propositional Logic). Then, because $t_2=\overline m$ is provable, $t_1\le t_2$ is provable by $A25$.
\section*{Problem 6}  
The part of this proof for the true atomic $\Sigma_0$ sentences was just Problem 5, so we only need to prove the part for the false sentences.\newline 
So let us consider a false $\Sigma_0$ sentence of the form $t_1=t_2$ where $t_1, t_2$ are both terms. Well, $\overline m\ne \overline n$ for the $m, n$ that 
$t_1, t_2$ designate respectively, is provable by $\Omega_3$. Thus, by two applications of $A25$ (one on $m$, the other on $n$), we can prove $\oldneg (t_1 =
t_2)$. In the other case, suppose that $t_1\le t_2$ is the false $\Sigma_0$ sentence. We choose $m, n$ similarly to the first part of this problem. Because 
$\overline m\le \overline n$ is false, $\overline m=0,...,\overline m=\overline n$ are all false, refutable by $\Omega_3$, and therefore, by Propositional 
Logic and $\Omega_4$, $\overline m\le \overline n$ is refutable. Then, by two applications of $A25$, $t_1\le t_2$ is refutable.  
\section*{Problem 7}
Assume that $F(\overline 0),F(\overline 1),...,F(\overline n)$ are all provable, and we are to show that the formula $(\forall x\le \overline n)F(x)$ is 
provable in PA. Recall that $(\forall x\le \overline n)F(x)$ is the same as $\forall x(x\le n\supset F(x)$. By $A25$, $F(\overline m)\supset ((x=\overline m
)\supset F(x)$ is provable for any number $n$. Hence, by M.P., we have the formulas $(x=\overline 0)\supset F(x),...,(x=\overline n)\supset F(x)$ all 
being provable, and thus, by Propositional logic, so is the formula $(x=\overline 0\vee ...\vee x=\overline n)\supset F(x)$. By $\Omega_4$ and Propositional 
logic, this implies $x\le \overline n\supset F(x)$, and thus, by Proposition 1(a), $\forall x(x\le n\supset F(x)$ is provable.\newline 
\section*{Problem 8}
We can prove that PA is $\Sigma_0$ complete by showing the every $\Sigma_0$ sentence is correctly decidable (meaning it defines itself), by using complete 
induction. In the base case, every atomic $\Sigma_0$ formula is correctly decidable by Problem 6. Let us make a complete induction assumption.
Consider a $\Sigma_0$ sentence of degree $n+1$. If it is of the form $X\wedge Y, X\vee Y, X\supset Y, \oldneg X$, 
$X$ and $Y$ are both correctly decidable, obviously meaning the sentence is also correctly decidable. 
\newline 
Next, we must consider the other two cases, namely, those $\Sigma_0$ sentences of the form 
$(\forall x\le \overline n)F(x)$ or $(\exists x\le \overline n)F(x)
$. In the first case, we assume that $(\forall x\le \overline n)F(x)$ 
is true. Now, this means that $F(\overline 0),...,F(\overline n)$ are all true, and thus,
 by the induction assumption, they are all correctly decidable, and hence provable, and thus, 
by Problem 7, so is $(\forall x\le \overline n)F(x)$. Now, assume that $(\forall x\le \overline n)F(x)$ is false. This implies that, for some $m\le n$, 
$F(\overline m)$ is false, 
and therefore, because $F(\overline m)$ is correctly decidable by the induction assumption, it is refutable. Now, $A10$ implies $\oldneg F(t)\supset 
\oldneg\forall xF(x)$. Well, we now know that $\overline m\le \overline n\wedge \oldneg F(\overline m)$ is provable (because all atomic $\Sigma_0$ 
formulas have been shown to be correctly decidable) and thus so is $\oldneg (\overline m\le \overline n\supset F(\overline m)$, and hence, so is
$\oldneg\forall x(x\le \overline n\supset F(x)$, which is the same as saying $\oldneg (\forall x\le \overline n)F(x)$, and thus we have shown 
$(\forall x\le \overline n)F(x)$ to be correctly decidable.\newline
On the other hand, consider $(\exists x\le \overline n)F(x)$. Assume that it is true. Hence, for some $m\le n$, $F(\overline m)$ is provable. 
Because $\overline m\le \overline n$ is also provable, so is $\overline m\le \overline n\supset F(\overline m)$, and thus, by $A11$, so is 
$\exists x(\overline x\le \overline n\supset F(x)$, which is the same as saying $(\exists x\le \overline n)F(x)$. Now, assume that $(\exists x\le \overline n
)F(x)$ is false. Thus, by the induction assumption, $\oldneg F(\overline 0),...,\oldneg F(\overline n)$ are all correctly decidable, 
true, and hence provable. 
This implies (by Problem 7) the sentence $(\forall x\le \overline n)\oldneg F(x)$ which implies $\forall x(x\le \overline n\supset \oldneg F(x)$, and thus, 
by $A10$, $x\le \overline n\supset \oldneg F(x)$ is provable, which implies $\oldneg(x\le \overline n\wedge F(x))$, and thus, by Proposition 1(a), 
so is $\forall x\oldneg(x\le \overline n\wedge F(x))$, and thus, by $A26$, $\oldneg\exists x(x\le \overline n\wedge F(x)$ is provable, which is the same 
as saying $\oldneg(\exists x\le \overline n)F(x)$.
\section*{Problem 9}
$L_1$: This is immediate from Problem 7.\newline
$L_2$: This is immediate from $\Omega_5$ and Proposition 1(a).
\end{document}

